% Nouveau EGA V (cote 25): archive cote 25, site page 33 (PDF page 34).
% typescript. Draft compared: none.
% AI-transcribed by Claude Opus 5.5 (claude-opus-5-5, adaptive thinking, effort high), Message Batches API, batch ega25b (msgbatch_01A7n5YuXubUPNK1zsepi8RG), 2026-11-13; not yet checked by eye.
\scanpage{25/33}
\mspage{29}

Noter qu'avec les notations précédentes, si $\hand{\zeta\in}$ $\check{P}$ - Z , alors pour tout
y $\hand{\in}$ Y$_{\hand{\zeta}}$ , on a coprof$_y$ F$_{k(\hand{\zeta})}$ = coprof$_y$ G$_{\hand{\zeta}}$ , et par suite: si coprof F$\hand{\leq}$n,
alors pour $\hand{\zeta\in}\check{P}$ - Z, on a coprof G$_{\hand{\zeta}}$ $\hand{\leq}$ n; en particulier si F est CM, alors
pour $\hand{\zeta\in}$ $\check{P}$ - Z, G$_{\hand{\zeta}}$ est CM. \inserted{Enfin, si F est $(S_k)$, \unsure{on sait que} $G_\zeta$ $(S_{k-1})$ pour $\zeta\in\check{P}-Z$ (réf EGA IV).}
\hand{8.3.} On notera que \deleted{on n'a pas de résultat analogue pour la propriété} \inserted{\illegible{6}}
\deleted{(S$_k$)}, même si F=$\underline{O}_X$ , k=1, X étant géom intègre de dim.2 (ou \inserted{k=2,} X étant
géométriquement intègre et géométriquement normal de dim.3). Il suffit
de partir d'une surface projective \xout{xxxxxx} intègre \inserted[typed]{X $\subset$ P$^r$} sur k algébriquement
clos, ayant un point \inserted[typed]{x} où X n'est pas Cohen-Macaulay; alors pour tout
hyperplan passant par x, la section hyperplane Y$_{\hand{\zeta}}$ correspondante admet
x pour cycle associé immergé (resp. on part d'une variété normale\hand{,} \hand{(}donc
\inserted[typed]{(S$_2$)\hand{)},} intègre X $\subset$ P$^r$ de dimension 3, ayant un point x $\in$ X en lequel X n'est pas
Cohen-Maaaulay, alors les Y passant par x ne sont pas \deleted{\unsure{(S}} CM i.e. ne sont
pas (S$_2$) en x)\hand{.} \over{d}{D}ans ces exemples, l'ensemble des \hand{$\zeta$} ``exceptionnels''
pour la propriété (S$_k$) contieht l'hyperplan d e $\check{P}$ défini par x$\in$P, et est
de codimension 1 (et non de codimension$\geq$2). \inserted{Cf \deleted{\illegible{1}} 8.5. plus bas pour \deleted{\unsure{un}} résultat général \unsure{précis} dans cette voie.}

\deleted{Cependant, si F satisfait à la propriété (S$_k$), on peut \unsure{montrer|monteer} à
l'aide de ce qui précède que l'ensemble des $\hand{\zeta\in}$ $\check{P}$ tels que $\phi_{\hand{\zeta}}$ soit
F$_{k(\hand{\zeta})}$-régulier\hand{,} \uline{et} G$_{\hand{\zeta}}$ satisfai\xout{\illegible{1}}t (S$_{k-1}$), a un complémentaire qui est de
codimmnsion $\geq$ 2. Cela résulte du critère 5.7.4. via la coppofondeur,
du fait que l'ensemble Z$_n$ des points de X où F est d e coprofondeur $>$n est
fermé, et que la codimension de Z$_{n\hand{\zeta}}$ dans X$_{\hand{\zeta}}$= Y$_{\hand{\zeta}}$ , même si elle peut
faire (pour d es \hand{$\zeta$} exceptionnels) des sauts \deleted{en hauteur} \inserted{vers le bas}, ledit saut ne peut
jamais être que d'une unité au plus, pour des raisons de dimension éviden-%
tes.}

\inserted{Si F est $(S_k)$, \unsure{alors} l'ens. des $\zeta\in\check{P}$ tels que $\phi_\zeta$ soit $F_{k(\zeta)}$-régulier et $G_\zeta$ soit $(S_k)$ ait un complémentaire de codim $\geq 2$,}

\margin{33}
